% =========================================================================
%  Dokumentacja techniczna: Stanowisko EGEA do diagnostyki zawieszenia.
%  Wersja 4.0 / czerwiec 2026  --- pełne wyjaśnienie krok po kroku.
%  Autor: aplikacja zgodna ze specyfikacją EGEA SPECSUS2018
% =========================================================================
\documentclass[11pt,a4paper]{article}

% ---- Pakiety podstawowe -------------------------------------------------
\usepackage[utf8]{inputenc}
\usepackage[T1]{fontenc}
\usepackage[polish]{babel}
\usepackage[a4paper, margin=2.4cm]{geometry}
\usepackage{lmodern}
\usepackage{microtype}

% ---- Matematyka ---------------------------------------------------------
\usepackage{amsmath, amssymb, amsthm, mathtools}
\usepackage{bm}

% ---- Grafika i tabele ---------------------------------------------------
\usepackage{graphicx}
\usepackage{booktabs}
\usepackage{array}
\usepackage{longtable}
\usepackage{multirow}
\usepackage{tabularx}
\usepackage{xcolor}
\usepackage{colortbl}

% ---- Pudełka informacyjne / wyjaśnienia ---------------------------------
\usepackage[most]{tcolorbox}
\tcbset{
  noteboxstyle/.style={
    colback=blue!4!white, colframe=blue!50!black,
    boxrule=0.4pt, arc=2pt, left=6pt, right=6pt, top=4pt, bottom=4pt,
    fonttitle=\bfseries,
  },
  intuitionstyle/.style={
    colback=yellow!8!white, colframe=orange!70!black,
    boxrule=0.4pt, arc=2pt, left=6pt, right=6pt, top=4pt, bottom=4pt,
    fonttitle=\bfseries,
  },
  examplestyle/.style={
    colback=green!4!white, colframe=green!40!black,
    boxrule=0.4pt, arc=2pt, left=6pt, right=6pt, top=4pt, bottom=4pt,
    fonttitle=\bfseries,
  },
  warningstyle/.style={
    colback=red!4!white, colframe=red!60!black,
    boxrule=0.4pt, arc=2pt, left=6pt, right=6pt, top=4pt, bottom=4pt,
    fonttitle=\bfseries,
  },
}
\newtcolorbox{notatka}[1][]{noteboxstyle, title=Definicja / Notatka, #1}
\newtcolorbox{intuicja}[1][]{intuitionstyle, title=Intuicja / Co to znaczy?, #1}
\newtcolorbox{przyklad}[1][]{examplestyle, title=Przykład z liczbami, #1}
\newtcolorbox{uwagabox}[1][]{warningstyle, title=Uwaga praktyczna, #1}

% ---- Hiperłącza i odsyłacze --------------------------------------------
\usepackage[hidelinks, breaklinks]{hyperref}
\usepackage{cleveref}

% ---- Listingi i kod -----------------------------------------------------
\usepackage{listings}
\definecolor{codekw}{RGB}{0,80,150}
\definecolor{codestr}{RGB}{160,40,40}
\definecolor{codecom}{RGB}{120,120,120}
\definecolor{codebg}{RGB}{248,248,248}
\lstdefinestyle{py}{
  language=Python,
  basicstyle=\ttfamily\footnotesize,
  keywordstyle=\color{codekw}\bfseries,
  stringstyle=\color{codestr},
  commentstyle=\color{codecom}\itshape,
  backgroundcolor=\color{codebg},
  numbers=left,
  numberstyle=\tiny\color{codecom},
  frame=single,
  rulecolor=\color{codecom!50},
  breaklines=true,
  showstringspaces=false,
  captionpos=b,
  tabsize=2,
  inputencoding=utf8,
  extendedchars=true,
  literate=%
    {ą}{{\k{a}}}1 {ć}{{\'c}}1 {ę}{{\k{e}}}1 {ł}{{\l{}}}1 {ń}{{\'n}}1
    {ó}{{\'o}}1 {ś}{{\'s}}1 {ź}{{\'z}}1 {ż}{{\.z}}1
    {Ą}{{\k{A}}}1 {Ć}{{\'C}}1 {Ę}{{\k{E}}}1 {Ł}{{\L{}}}1 {Ń}{{\'N}}1
    {Ó}{{\'O}}1 {Ś}{{\'S}}1 {Ź}{{\'Z}}1 {Ż}{{\.Z}}1,
}

% ---- Layout tytułu ------------------------------------------------------
\usepackage{titlesec}
\titleformat{\section}
  {\normalfont\Large\bfseries\color{codekw}}{\thesection}{1em}{}
\titleformat{\subsection}
  {\normalfont\large\bfseries\color{codekw!80!black}}{\thesubsection}{1em}{}

% ---- Definicje matematyczne --------------------------------------------
\newcommand{\R}{\mathbb{R}}
\newcommand{\dd}{\,\mathrm{d}}
\newcommand{\diff}[2]{\dfrac{\mathrm{d}#1}{\mathrm{d}#2}}
\newcommand{\ddtwo}[2]{\dfrac{\mathrm{d}^2#1}{\mathrm{d}#2^2}}
\renewcommand{\vec}[1]{\mathbf{#1}}
\DeclareMathOperator{\arctantwo}{arctan2}

\newtheoremstyle{plainpol}{\topsep}{\topsep}{\itshape}{}%
  {\bfseries}{.}{0.5em}{\thmname{#1}\thmnumber{ #2}\thmnote{ (\itshape #3\upshape)}}
\theoremstyle{plainpol}
\newtheorem{definicja}{Definicja}[section]
\newtheorem{twierdzenie}[definicja]{Twierdzenie}
\newtheorem{uwaga}[definicja]{Uwaga}

\title{\textbf{Stanowisko EGEA do diagnostyki zawieszenia pojazdu}\\[0.4em]
  \large Wyjaśnienie krok po kroku --- od fizyki przez matematykę do kodu \\
  metoda minimalnego przesunięcia fazowego $\varphi_{\min}$ wg normy EGEA SPECSUS2018}
\author{Aplikacja v4.0 -- model 2 DOF + 1 DOF, dokumentacja pedagogiczna}
\date{czerwiec 2026}

% =========================================================================
\begin{document}
% =========================================================================

\maketitle
\tableofcontents
\clearpage

% =========================================================================
\section{Wprowadzenie dla początkujących}
% =========================================================================

\subsection{O czym jest ten dokument?}

Ten dokument tłumaczy, \textbf{jak działa stanowisko diagnostyczne}, na którym
sprawdza się stan amortyzatorów w samochodzie -- nie demontując ich z auta.
Wyjaśnienia są pisane bardzo prostym językiem, z dużą liczbą rysunków
pojęciowych i przykładów liczbowych. Każdy wzór ma wytłumaczenie, co
fizycznie znaczą występujące w nim symbole i co się stanie z wynikiem,
jeśli jeden z nich będzie większy lub mniejszy.

\begin{intuicja}
\textbf{Najprostszy obraz całości:} samochód najeżdża kołem na płytę, która
\emph{trzęsie się} w pionie. Ta wibracja przechodzi przez oponę i przez
amortyzator do nadwozia. Stanowisko mierzy siłę nacisku koła na płytę
w czasie, a komputer porównuje rytm tej siły z rytmem ruchu płyty.
\textbf{Jeżeli siła nadąża dokładnie za płytą} -- amortyzator jest zużyty
(brak tłumienia). \textbf{Jeżeli siła wyraźnie spóźnia się względem
płyty} -- amortyzator dobrze rozprasza energię i jest sprawny. Tę różnicę
w czasie mierzymy jako \emph{kąt fazowy} $\varphi$.
\end{intuicja}

\subsection{Trzy najważniejsze pojęcia (bez wzorów)}

\begin{enumerate}
  \item \textbf{Tłumik / amortyzator $(c_M)$.} To element zawieszenia, który
    \emph{rozprasza energię drgań} w postaci ciepła. Pomyśl o nim jak
    o \emph{strzykawce wypełnionej oliwą}: kiedy próbujesz ją szybko ścisnąć,
    odczuwasz opór. Im szybciej ściskasz, tym większy opór. Sprawny tłumik
    sprawia, że auto po najechaniu na dziurę raz się ugnie i wraca na
    miejsce. Niesprawny -- po dziurze auto huśta się jak na sprężynie.

  \item \textbf{Sprężyna $(k_M, k_m)$.} Element gromadzący energię potencjalną.
    Sprężyna sama nie rozprasza energii -- gdy ją ściśniesz, oddaje
    całą energię z powrotem. W aucie sprężyna utrzymuje nadwozie nad
    kołem, ale to tłumik decyduje, jak szybko drganie wygaśnie.

  \item \textbf{Faza / przesunięcie fazowe $(\varphi)$.} Mówi, \emph{jak
    bardzo dwa rytmiczne zjawiska są ze sobą zsynchronizowane}. Wyobraź
    sobie dwóch tancerzy: jeśli kroczą w tym samym momencie -- faza
    $0^\circ$. Jeśli drugi tancerz spóźnia się o pół taktu -- faza
    $180^\circ$. W naszym przypadku ,,tancerze'' to: ruch płyty
    stanowiska i siła nacisku koła. Im więcej spóźnia się siła, tym
    bardziej amortyzator tłumi (bo opór tłumika rośnie z prędkością, a nie
    z położeniem).
\end{enumerate}

\subsection{Dlaczego ,,minimalna'' faza, a nie po prostu faza?}

Stanowisko \emph{przelatuje przez całe spektrum częstotliwości} od 18 Hz
w dół do 6 Hz. W każdym momencie tej przeciągłej zjazdówki sprawdzamy,
\emph{ile wynosi $\varphi$}. Wartość, w której $\varphi$ osiąga
\textbf{najmniejszą} wartość w całym pomiarze, nazywamy
$\varphi_{\min}$ i ona jest miarą stanu amortyzatora. Zasada decyzyjna
EGEA SPECSUS2018:
\begin{itemize}
  \item $\varphi_{\min} \geq 35^\circ \Rightarrow$ amortyzator
    \textbf{sprawny};
  \item $\varphi_{\min} < 35^\circ \Rightarrow$ amortyzator
    \textbf{niesprawny / podejrzany}.
\end{itemize}

\begin{uwagabox}
\textbf{Ograniczenie modelu ćwiartkowego (2 DOF).} W aplikacji uzupełniamy
to kryterium o dwa pomocnicze wskaźniki: \emph{EUSAMA} oraz wykrycie
,,odrywania koła'' ($F_{\min}<0$). Dla niektórych konfiguracji wartości
$k_m, c_m$ samego rezonansu koła wzbudza się tak mocno, że $\varphi_{\min}$
nie spada poniżej $35^\circ$ pomimo ewidentnej awarii. Dwa pomocnicze
kryteria zapobiegają fałszywej diagnozie ,,sprawny'' w takiej sytuacji
(zob. \cref{sec:diag_app}).
\end{uwagabox}

\subsection{Plan dokumentu}

\begin{enumerate}
  \item \textbf{Geneza metody} -- skąd się wzięła i co zastąpiła
    (\cref{sec:geneza}).
  \item \textbf{Fizyka} -- co modelujemy i dlaczego (\cref{sec:model}).
  \item \textbf{Matematyka równań ruchu} -- jak zapisujemy fizykę
    formalnie (\cref{sec:rownania}).
  \item \textbf{Procedura MinPhaseShift krok po kroku} -- algorytm
    obliczania $\varphi_{\min}$ z mierzonego sygnału (\cref{sec:procedura}).
  \item \textbf{Model 1 DOF} -- prostsza wersja do analizy parametrycznej
    (\cref{sec:1dof}).
  \item \textbf{Wskaźniki diagnostyczne} -- co aplikacja wypisuje
    (\cref{sec:wskazniki}).
  \item \textbf{Kryteria oceny} -- jak interpretować liczby
    (\cref{sec:kryteria}).
  \item \textbf{Implementacja w kodzie Python} (\cref{sec:kod}).
\end{enumerate}

% =========================================================================
\section{Geneza i kontekst metody}\label{sec:geneza}
% =========================================================================

\subsection{Co było wcześniej -- metoda EUSAMA}

Pierwsza znormalizowana metoda płytowego testowania amortyzatorów
powstała w 1976~r. (EUSAMA~\cite{EUSAMA1976}). Idea jest bardzo prosta:
płyta drga jednostajnie sinusoidalnie, mierzy się \emph{najmniejszą}
chwilową siłę nacisku $F_{\min}$ i dzieli przez ciężar statyczny $F_{st}$.
Wynikiem jest procent przylegania koła do nawierzchni:
\begin{equation}
  \mathrm{EUSAMA} = \dfrac{F_{\min}}{F_{st}}\cdot 100\,\%.
  \label{eq:eusama}
\end{equation}

\begin{intuicja}
Wyobraź sobie, że płyta podskakuje w górę i w dół. Im słabszy amortyzator,
tym mocniej koło ,,odbija się'' i \emph{traci kontakt z płytą} w pewnych
momentach. $F_{\min}$ to siła w tym najgorszym momencie. Jeśli koło
naprawdę się odrywa -- $F_{\min} = 0$ i EUSAMA = 0\%. Jeśli amortyzator
trzyma świetnie -- $F_{\min}\approx F_{st}$ i EUSAMA $\approx 100\%$.
\end{intuicja}

\subsection{Dlaczego EUSAMA przestała wystarczać?}

Klapka i inni~\cite{Klapka2017} pokazali eksperymentalnie, że wynik
EUSAMA zależy bardziej od:
\begin{itemize}
  \item masy badanego pojazdu,
  \item ciśnienia w oponach,
  \item geometrii zawieszenia,
\end{itemize}
niż od \emph{rzeczywistego stanu amortyzatora}. Czyli ta sama wartość
EUSAMA może oznaczać sprawny amortyzator w aucie A i zużyty w aucie B
-- co czyni metodę mało wiarygodną w warunkach STKP.

\subsection{Metoda fazowa Tsymberowa (1996)}

Tsymberow~\cite{Tsymberov1996} zaproponował zupełnie inną filozofię:
nie patrzymy na \emph{wartość} siły, lecz na \emph{moment}, w którym
osiąga ona swoje ekstrema względem ruchu płyty. Ta różnica czasowa
jest miarą \emph{przesunięcia fazowego} i (jak pokazano teoretycznie)
znacznie słabiej zależy od masy pojazdu.

\begin{intuicja}
Jeżeli amortyzator jest niesprawny, układ koło--zawieszenie zachowuje się
,,jak sprężyna'' -- prawie bez dyssypacji energii. Siła reakcji i ruch
płyty są w fazie (albo w przeciwfazie). Jeśli amortyzator pracuje
prawidłowo, dyssypuje energię, a faza siły wyraźnie spóźnia się
o kilkadziesiąt stopni względem ruchu płyty. \emph{Im większe to
spóźnienie, tym lepszy amortyzator.}
\end{intuicja}

\subsection{Standaryzacja EGEA SPECSUS2018}

European Garage Equipment Association uznała metodę MPS za europejski
standard~\cite{EGEA2019}. Zdefiniowano: profil zmian częstotliwości,
sposób detekcji $\varphi(i)$, próg kwalifikacji $\varphi_{\min} =
35^\circ$ i wszystkie dopuszczalne tolerancje.

% =========================================================================
\section{Model fizyczny stanowiska}\label{sec:model}
% =========================================================================

\subsection{Schemat układu (rysunek pojęciowy)}

\begin{center}
\fbox{\begin{tabular}{c}
    \texttt{[ NADWOZIE -- masa $M$ ]}\\[0.3em]
    \texttt{|\,||\,|} \quad sprężyna $k_M$ \emph{równolegle} z tłumikiem $c_M$ \\[0.3em]
    \texttt{[ KOŁO + PIASTA -- masa $m$ ]}\\[0.3em]
    \texttt{|\,||\,|} \quad sprężyna $k_m$ \emph{równolegle} z tłumikiem $c_m$ \quad (opona) \\[0.3em]
    \texttt{[ PŁYTA STANOWISKA, $z(t) = d\cos\theta(t)$ ]}
\end{tabular}}
\end{center}

\begin{notatka}[title=Co modelujemy]
Tylko \textbf{jedno koło} i odpowiadającą mu ćwiartkę nadwozia -- stąd
nazwa \emph{model ćwiartkowy} (ang.~\emph{quarter car}). Dwa stopnie
swobody (2 DOF) to:
\begin{itemize}
  \item $x_m$ -- pionowe przemieszczenie koła,
  \item $x_M$ -- pionowe przemieszczenie nadwozia.
\end{itemize}
Wszystkie pozostałe wielkości (sztywności, tłumienia, masy, amplituda
płyty) są \emph{parametrami}, czyli liczbami stałymi w czasie symulacji.
\end{notatka}

\subsection{Co reprezentuje każdy z parametrów?}

\begin{center}
\begin{tabular}{lp{8.5cm}l}
\toprule
\textbf{Symbol} & \textbf{Znaczenie fizyczne} & \textbf{Typowa wartość} \\
\midrule
$M$ & masa nadwozia przypadająca na jedno koło & 346~kg \\
$m$ & masa nieresorowana (koło, hamulec, dolne ramię) & 36~kg \\
$k_M$ & sztywność \emph{sprężyny} zawieszenia (resoru, sprężyny śrubowej) & 25\,570~N/m \\
$c_M$ & tłumienie amortyzatora -- to chcemy diagnozować! & 1\,474~N$\cdot$s/m \\
$k_m$ & sztywność opony (jak sprężyna pod kołem) & 253\,161~N/m \\
$c_m$ & tłumienie wewnętrzne opony (małe, ale niezerowe) & 150~N$\cdot$s/m \\
$d$ & amplituda drgania płyty (połowa peak-to-peak) & 0,003~m = 3~mm \\
$m_p$ & masa samej płyty (potrzebna do kalibracji) & 20~kg \\
\bottomrule
\end{tabular}
\end{center}

\begin{intuicja}
Zwróć uwagę: \textbf{opona jest o rząd wielkości sztywniejsza} niż
sprężyna zawieszenia ($k_m / k_M \approx 10$). To znaczy, że dla małych
amplitud opona zachowuje się prawie jak ,,sztywny dystans'' -- a ruch
koła $x_m$ podąża blisko ruchu płyty $z$. Natomiast \emph{tłumienie}
opony $c_m$ jest dziesiątki razy mniejsze niż $c_M$ -- dlatego za
rozpraszanie energii w pasmie 6--18~Hz odpowiada głównie amortyzator.
\end{intuicja}

\subsection{Wymuszenie -- ruch płyty}

Płyta jest napędzana przez wał z mimośrodem, więc porusza się jak
\emph{punkt na okręgu rzutowany na oś pionową}:
\begin{equation}
  z(t) = d\cos(\theta(t)), \qquad
  \theta(t) = \int_0^t 2\pi f(\tau)\,d\tau.
  \label{eq:z}
\end{equation}

\begin{intuicja}
Funkcja $\cos$ zaczyna się od wartości 1 (płyta na samej górze --
\emph{TOP}), opada przez 0 do $-1$ (płyta na dole --
\emph{BOTTOM}), wraca przez 0 do 1 (kolejne TOP). Argument $\theta(t)$
to \emph{aktualna faza}: $\theta=0$ to TOP, $\theta=\pi/2$ to przejście
przez środek w drodze w dół, $\theta=\pi$ to BOTTOM, $\theta=3\pi/2$ to
przejście przez środek w drodze w górę, $\theta=2\pi$ wracamy do TOP.
Wzór z całką oznacza po prostu: ,,faza rośnie tym szybciej, im większa
częstotliwość'' -- jeśli $f$ jest stałe, to $\theta(t) = 2\pi f\cdot t$,
czyli klasyczny ruch harmoniczny.
\end{intuicja}

% =========================================================================
\section{Równania ruchu -- skąd się biorą}\label{sec:rownania}
% =========================================================================

\subsection{II zasada Newtona, krok po kroku}

Dla każdej masy stosujemy $\sum F = m\cdot a$. Zacznijmy od koła
(masy nieresorowanej $m$). Co na nie działa?

\begin{enumerate}
  \item \textbf{Siła od sprężyny zawieszenia $k_M$:} sprężyna jest
    rozciągnięta/ściśnięta o $(x_m - x_M)$ -- różnica wysokości koła
    i nadwozia. Siła to $-k_M(x_m - x_M)$ (minus, bo sprężyna ciągnie
    w kierunku redukcji rozciągnięcia).

  \item \textbf{Siła od tłumika zawieszenia $c_M$:} tłumik daje siłę
    proporcjonalną do prędkości względnej (jak strzykawka z oliwą):
    $-c_M(\dot{x}_m - \dot{x}_M)$.

  \item \textbf{Siła od sprężyny opony $k_m$:} opona jest rozciągnięta
    o $(x_m - z)$ -- różnica położenia koła i płyty:
    $-k_m(x_m - z)$.

  \item \textbf{Siła od tłumienia opony $c_m$:} $-c_m(\dot{x}_m - \dot{z})$.
\end{enumerate}

Sumując wszystkie cztery siły i przyrównując do $m\cdot a = m\ddot{x}_m$:
\begin{equation}
  m\,\ddot{x}_m = -k_M(x_m - x_M) - c_M(\dot{x}_m - \dot{x}_M)
                  - k_m(x_m - z) - c_m(\dot{x}_m - \dot{z}).
  \label{eq:m_raw}
\end{equation}

\begin{intuicja}
Każdy minus oznacza, że dany element ,,ciągnie z powrotem''. Jeśli koło
podskoczy w górę ($x_m > 0$), to sprężyny próbują je ściągnąć w dół.
Jeśli porusza się w górę ($\dot{x}_m > 0$), to tłumiki hamują ten ruch.
W ten sposób cała mechanika sprzęga się z \emph{ruchem płyty} $z(t)$
poprzez ostatnie dwa wyrazy.
\end{intuicja}

Dla masy resorowanej $M$ (nadwozia) działają tylko siły od zawieszenia,
nie od opony:
\begin{equation}
  M\,\ddot{x}_M = +k_M(x_m - x_M) + c_M(\dot{x}_m - \dot{x}_M).
  \label{eq:M_raw}
\end{equation}

Zwróć uwagę na znak: tu jest \emph{plus}, bo gdy koło $x_m$ wzrasta
względem nadwozia, sprężyna popycha nadwozie do góry.

\subsection{Uporządkowanie -- standardowa forma}

Rozwijamy nawiasy w równaniach \cref{eq:m_raw,eq:M_raw} i porządkujemy
po niewiadomych:
\begin{equation}
\boxed{\begin{aligned}
  m\,\ddot{x}_m &= -(k_M+k_m)x_m + k_Mx_M -(c_M+c_m)\dot{x}_m + c_M\dot{x}_M + k_m z + c_m \dot{z}, \\
  M\,\ddot{x}_M &= k_Mx_m - k_Mx_M + c_M\dot{x}_m - c_M\dot{x}_M.
\end{aligned}}
\label{eq:uklad}
\end{equation}

\begin{intuicja}
Po lewej mamy ,,bezwładność'' (masa razy przyspieszenie). Po prawej
\emph{sumę wszystkich sił} podzieloną na cztery grupy:
\begin{itemize}
  \item co zależy od położenia koła $x_m$,
  \item co zależy od położenia nadwozia $x_M$,
  \item co zależy od prędkości koła $\dot{x}_m$,
  \item co zależy od prędkości nadwozia $\dot{x}_M$.
\end{itemize}
Te dwie ostatnie kolumny ($k_m z + c_m \dot{z}$) to ,,wejście''
sterujące -- ruch płyty, na który nie mamy wpływu (to wymuszenie).
\end{intuicja}

\subsection{Postać macierzowa -- po co to robimy?}

Dwa równania drugiego rzędu zamieniamy na cztery równania pierwszego
rzędu, wprowadzając \emph{wektor stanu}:
\begin{equation}
  \vec{X}(t) = \begin{bmatrix} x_m(t) \\ x_M(t) \\ \dot{x}_m(t) \\ \dot{x}_M(t) \end{bmatrix} \in \R^4.
\end{equation}

\begin{intuicja}
Wektor stanu to po prostu \emph{lista czterech liczb}, które w pełni
opisują, ,,co się dzieje'' w danej chwili. Pierwsze dwie -- gdzie są
masy, dwie kolejne -- jak szybko się poruszają. Mając tę listę
w chwili $t$ i znając równania ruchu, możemy obliczyć, gdzie będzie
ona za moment ($t + dt$). To pozwala wprost zastosować numeryczną
metodę całkowania (\texttt{scipy.integrate.odeint}).
\end{intuicja}

\begin{twierdzenie}[Postać macierzowa]
Układ \cref{eq:uklad} jest równoważny liniowemu układowi pierwszego rzędu:
\begin{equation}
  \dot{\vec{X}}(t) = \vec{A}\vec{X}(t) + \vec{B}(t),
  \label{eq:stan}
\end{equation}
gdzie:
\begin{equation}
  \vec{A} = \begin{bmatrix}
    0 & 0 & 1 & 0 \\
    0 & 0 & 0 & 1 \\
    -\dfrac{k_m+k_M}{m} & \dfrac{k_M}{m} & -\dfrac{c_m+c_M}{m} & \dfrac{c_M}{m} \\[6pt]
    \dfrac{k_M}{M} & -\dfrac{k_M}{M} & \dfrac{c_M}{M} & -\dfrac{c_M}{M}
  \end{bmatrix},\qquad
  \vec{B}(t) = \begin{bmatrix} 0 \\ 0 \\ \dfrac{k_m z(t) + c_m \dot{z}(t)}{m} \\ 0 \end{bmatrix}.
\end{equation}
\end{twierdzenie}

\begin{intuicja}
Pierwsze dwa wiersze macierzy $\vec{A}$ to po prostu zapis: ,,pochodna
położenia to prędkość'' (definicja). Trzeci wiersz to równanie dla koła
po podzieleniu obu stron przez $m$. Czwarty -- analogicznie dla nadwozia.
Wektor $\vec{B}(t)$ zawiera tylko to, co \emph{wpływa z zewnątrz} -- czyli
$z(t)$ i $\dot{z}(t)$ skalowane sztywnością i tłumieniem opony.
\end{intuicja}

\subsection{Wartości własne -- gdzie są rezonanse?}

Dla domyślnych parametrów numerycznie liczymy
$\det(\vec{A} - \lambda \vec{I}) = 0$ i otrzymujemy cztery wartości
zespolone $\lambda_k = -\zeta_k\omega_k \pm i\omega_k\sqrt{1-\zeta_k^2}$.
Po wyciągnięciu części urojonych i przeliczeniu na hertze:
\begin{align}
  f_1 &\approx 1{,}4~\mathrm{Hz} \quad \text{(rezonans nadwozia)},\\
  f_2 &\approx 13{,}5~\mathrm{Hz} \quad \text{(rezonans koła)}.
\end{align}

\begin{intuicja}
Każda swobodna masa drga przy określonej częstotliwości, podobnie jak
struna basowa drga przy innej częstotliwości niż wysoka. Nadwozie
(ciężkie, na miękkiej sprężynie) ma niską częstotliwość $f_1$. Koło
(lekkie, na sztywnej oponie) ma wysoką $f_2$. \emph{Norma celowo}
wprowadza zakres pomiarowy 6--18~Hz, by \textbf{zawrzeć w nim rezonans
koła}, bo to wokół niego dynamika układu jest najbardziej czuła na
zmiany tłumienia $c_M$.
\end{intuicja}

% =========================================================================
\section{Siła kontaktu i sygnały pomiarowe}
% =========================================================================

\subsection{Siła kontaktu opony $F(t)$}

\begin{definicja}[$F(t)$, rozdz.~3.6 normy]
Pionowa siła z jaką koło wciska oponę w płytę:
\begin{equation}
  F(t) = F_{st} + k_m\bigl(x_m - z\bigr) + c_m\bigl(\dot{x}_m - \dot{z}\bigr).
  \label{eq:F}
\end{equation}
gdzie $F_{st} = (M+m)g$ to ciężar przypadający na to koło.
\end{definicja}

\begin{intuicja}
$F$ rozkładamy na trzy człony:
\begin{enumerate}
  \item $F_{st}$ -- \textbf{stała} siła ciężkości (ciężar auta).
  \item $k_m(x_m - z)$ -- \textbf{sprężynowa} część siły kontaktu (im
    bardziej opona jest ściśnięta, tym mocniej napiera na płytę).
  \item $c_m(\dot{x}_m - \dot{z})$ -- \textbf{tłumieniowa} część siły
    (zwykle mała, ale dodaje fazę).
\end{enumerate}
\textbf{Jeśli $F(t)$ spadnie do zera lub poniżej} -- koło chwilowo
unosi się nad płytą. To sygnał poważnego problemu z amortyzatorem.
\end{intuicja}

\subsection{Siła pustej płyty $F_p(t)$ (do kalibracji)}

\begin{definicja}[$F_p(t)$, rozdz.~3.4 normy]
Sama płyta o masie $m_p$ też jest bezwładna i wymaga siły, żeby ją
przyspieszyć:
\begin{equation}
  F_p(t) = m_p\,\ddot{z}(t).
  \label{eq:Fp}
\end{equation}
\end{definicja}

\begin{intuicja}
W teście bez pojazdu czujnik siły pod płytą mierzy właśnie $F_p$. To
nasz \emph{zegar wzorcowy} -- z $F_p$ wiemy dokładnie, kiedy płyta jest
na samej górze (TOP). To pozwoli nam później precyzyjnie liczyć fazę,
gdy w teście z pojazdem czujnik mierzy sumę $F + F_p$.
\end{intuicja}

\subsection{Surowy sygnał $F_r(t)$ z tensometrów}

\begin{equation}
  F_r(t) = F(t) + F_p(t) \quad \Longrightarrow \quad F(t) = F_r(t) - F_p(t).
\end{equation}

\begin{notatka}[title=Po co kalibrować?]
Bez znajomości $F_p$ nie wiemy, ile siły pochodzi od opony, a ile od
samej trzęsącej się płyty. Dlatego \emph{przed} pierwszym pomiarem
robimy test bez pojazdu, mierząc $F_p(t)$. Macierz tych pomiarów
zapamiętujemy raz; potem używamy jej przy każdym kolejnym aucie.
\end{notatka}

\subsection{Sygnał wyzwalania $S(t)$ i znaczniki $ST_i$}

\begin{definicja}[$ST_i$, rozdz.~3.8 normy]
Z czujnika indukcyjnego pod płytą dostajemy sygnał 0/1: $S(t)=1$ w
okolicach skrajnego (umownie -- najwyższego) położenia płyty. Znaczniki
$ST_i$ to \emph{pierwsze} próbki o wartości 1 w każdej grupie sąsiednich
jedynek:
\begin{equation}
  ST_i = \min\{k\in\mathbb{N}\colon S_k = 1 \wedge S_{k-1} = 0\}.
\end{equation}
\end{definicja}

\begin{intuicja}
$S$ to ,,migająca lampka'' zsynchronizowana z mechanizmem mimośrodu.
Każde mignięcie oznacza nowy cykl płyty. \textbf{Ale lampka nie zawsze
zapala się dokładnie w TOP} -- jest pewien luz/przesunięcie mechaniczne.
Dlatego wprowadzimy korekcję $\Delta\mathrm{Period}(i)$ liczoną z testu
bez pojazdu (\cref{sec:kalibracja}).
\end{intuicja}

% =========================================================================
\section{Funkcja zmienności częstotliwości}
% =========================================================================

\subsection{Pięć faz testu -- po co taki skomplikowany profil?}

Norma rozbija test na pięć faz, każda ma własne zadanie:

\begin{center}
\begin{tabular}{cllc}
\toprule
\textbf{Faza} & \textbf{Co robi} & \textbf{Czas trwania} & \textbf{$f$ [Hz]} \\
\midrule
1 & rozruch płyty (rozpędzanie)           & $T_1 = 2$~s                 & $2 \to 25$ \\
2 & stabilizacja (pomiar $H_{25}$)        & $\Delta T_{25} = F_{st}\cdot 0{,}16 + 1200$~[ms] & $25$ \\
3 & schodzenie do zakresu pomiarowego     & $T_{\text{prep}} = 2$~s     & $25 \to 18$ \\
4 & WŁAŚCIWY POMIAR ($\varphi_{\min}$)    & $\Delta T_{\text{meas}} = 7{,}5$~s & $18 \to 6$ \\
5 & wygaszenie                            & $T_{\text{ext}} = 3$~s      & $6 \to 0$ \\
\bottomrule
\end{tabular}
\end{center}

\begin{intuicja}
\textbf{Po co fazy 1, 2, 3, 5}? Bo gdyby od razu zacząć od fazy 4
i ostro zakończyć po niej, układ \emph{nie ustabilizowałby się} -- na
początku miałby silne drgania własne, a na końcu wibracje zamiast
łagodnie zaniknąć, raptem ucichłyby (ryzyko mechanicznego szarpnięcia).
Faza 2 daje dodatkowo pomiar $H_{25}$ przy stabilnej, wysokiej
częstotliwości -- z niego oszacujemy sztywność opony rig.
\end{intuicja}

\subsection{Wzór analityczny na $f(t)$}

Wprowadzamy oznaczenia: $T_2 = T_1 + \Delta T_{25}$,
$T_3 = T_2 + T_{\text{prep}}$, $T_4 = T_3 + \Delta T_{\text{meas}}$.
Funkcja zmiany częstotliwości:
\begin{equation}
  f(t) =
  \begin{cases}
    2 + \dfrac{f_{\max}-2}{T_1}\,t, & 0 \leq t < T_1 \\[4pt]
    f_{\max} = 25, & T_1 \leq t < T_2 \\[4pt]
    25 - \dfrac{25-18}{T_{\text{prep}}}(t - T_2), & T_2 \leq t < T_3 \\[4pt]
    18 - \dfrac{18-6}{\Delta T_{\text{meas}}}(t - T_3), & T_3 \leq t < T_4 \\[4pt]
    6 - \dfrac{6}{T_{\text{ext}}}(t - T_4), & T_4 \leq t < T_4 + T_{\text{ext}} \\[4pt]
    0, & \text{w przeciwnym razie.}
  \end{cases}
  \label{eq:f_target}
\end{equation}

\begin{intuicja}
Każdy wiersz to po prostu \emph{linia prosta} przechodząca z jednej
wartości częstotliwości do drugiej w określonym czasie (rampa).
Liczba $0{,}16$ przy $F_{st}$ to empiryczny współczynnik zapewniający,
że cięższe pojazdy dostaną dłuższą stabilizację (bardziej bezwładne).
\end{intuicja}

\subsection{Cykliczność -- dlaczego w obrębie cyklu częstotliwość jest stała?}

Płyta jest napędzana mimośrodem (kołem korbowym), więc każdy obrót wału
realizuje jeden cykl drgania -- i \emph{w obrębie tego obrotu}
częstotliwość jest praktycznie stała (wał obraca się jednostajnie w tym
czasie). Dopiero między kolejnymi cyklami silnik delikatnie zmienia
obroty zgodnie z \cref{eq:f_target}.

Numerycznie generujemy więc cykle krok po kroku:
\begin{equation}
  t_{k+1} = t_k + \dfrac{1}{f(t_k)},
\end{equation}
a wewnątrz każdego cyklu fazę wyliczamy proporcjonalnie:
\begin{equation}
  \theta(t) = 2\pi\cdot\dfrac{t - t_k}{t_{k+1} - t_k}\quad \text{dla}\quad t\in[t_k, t_{k+1}).
\end{equation}

\begin{intuicja}
$\theta$ rośnie liniowo od $0$ (w chwili $t_k$, czyli TOP) do $2\pi$
(w chwili $t_{k+1}$, koniec cyklu, znów TOP). Wewnątrz cyklu mamy
,,porcję'' standardowego sinusoidalnego ruchu.
\end{intuicja}

% =========================================================================
\section{Procedura MinPhaseShift -- algorytm krok po kroku}\label{sec:procedura}
% =========================================================================

\subsection{Dane wejściowe}

Niezależnie od tego, skąd wzięliśmy sygnały (z symulatora czy z czujników),
mamy do dyspozycji trzy zsynchronizowane wektory:
\begin{equation}
  \mathrm{Ramka} = \bigl\{(t_i, F_i, S_i)\bigr\}_{i=0}^{N-1},
\end{equation}
gdzie:
\begin{itemize}
  \item $t_i$ -- chwile czasu (równomiernie spróbkowane);
  \item $F_i$ -- siła kontaktu opona-płyta w chwili $t_i$;
  \item $S_i \in \{0, 1\}$ -- sygnał czujnika TOP płyty.
\end{itemize}

Plus znana wartość $F_{st}$ -- statyczne obciążenie koła.

\subsection{Krok 1. Wygładzenie sygnału $F(t)$}

W rzeczywistości czujniki rejestrują też szum mechaniczny, drgania
o wysokich częstotliwościach. Przepuszczamy $F$ przez \textbf{filtr
dolnoprzepustowy} Butterwortha 4-go rzędu o częstotliwości odcięcia
$f_c = 50$~Hz. \emph{Bardzo ważne:} stosujemy filtrowanie dwukierunkowe
(\texttt{scipy.signal.filtfilt}), które \textbf{nie wprowadza
przesunięcia fazowego} -- to konieczne, bo cała metoda opiera się na
precyzyjnym mierzeniu fazy.

\begin{intuicja}
Wyobraź sobie zaszumione zdjęcie. Filtrowanie to ,,uśrednianie sąsiednich
pikseli'' -- usuwa szum, ale może też ,,rozmazać'' krawędzie. Filtr
\texttt{filtfilt} robi to dwa razy (raz w przód, raz w tył), dzięki czemu
rozmycia z obu kierunków znoszą się i krawędzie nie przesuwają się
w czasie.
\end{intuicja}

\subsection{Krok 2. Wykrywanie znaczników $ST_i$}

\begin{equation}
  \{ST_i\}_{i=1,\ldots,N_p} = \{k\colon S_k = 1 \wedge S_{k-1} = 0\}.
\end{equation}

Czyli przechodzimy przez $S$ i zaznaczamy momenty, w których sygnał
,,zapala się'' (przechodzi z 0 na 1). Każde takie ,,zapalenie'' to
początek nowego cyklu płyty.

Z dwóch kolejnych znaczników obliczamy \emph{okres} cyklu:
\begin{equation}
  \mathrm{Period}(i) = t_{ST_{i+1}} - t_{ST_i}, \qquad f(i) = \dfrac{1}{\mathrm{Period}(i)}.
\end{equation}

Wszystko, co dalej, robimy \emph{osobno dla każdego cyklu $i$}.

\subsection{Krok 3. Kalibracja dynamiczna -- macierz $\Delta\mathrm{Period}(i)$}\label{sec:kalibracja}

\begin{definicja}[$\Delta\mathrm{Period}(i)$, rozdz.~3.9 i~3.10 normy]
Dla testu \emph{bez pojazdu} i dla każdego cyklu $i$ wyznaczamy:
\begin{align}
  \mathrm{CalcTOP}(i) &= \arg\min_{t\in[ST_i,\,ST_{i+1}]} F_p(t), \\
  \Delta\mathrm{Period}(i) &= \mathrm{CalcTOP}(i) - ST_i.
\end{align}
\end{definicja}

\begin{intuicja}
\textbf{Po co \emph{minimum} $F_p$, a nie maksimum?} Pamiętajmy, że
$F_p = m_p\ddot{z}$, a $z = d\cos\theta$, więc $\ddot{z} = -d\omega^2\cos\theta$.
Najbardziej \emph{ujemne} (czyli minimum) $\ddot{z}$ wystąpi przy
$\cos\theta = 1$, czyli w \textbf{TOP płyty}. Argmin wybiera właśnie tę
chwilę. Liczba $\Delta\mathrm{Period}(i)$ mówi: \emph{,,jak bardzo
sygnał zapalający się $ST_i$ wyprzedza prawdziwe TOP płyty''}.
\end{intuicja}

\subsection{Krok 4. $TOP_p(i)$ dla testu z pojazdem}

Po skalibrowaniu już wiemy, jak ,,wyrównać'' nieprecyzyjny czujnik.
W teście z pojazdem:
\begin{equation}
  TOP_p(i) = ST(i) + \Delta\mathrm{Period}(i).
  \label{eq:TOPp}
\end{equation}

Czyli: do nieprecyzyjnego $ST(i)$ dodajemy znaną już korektę -- otrzymujemy
\emph{prawdziwy moment najwyższego położenia płyty} w cyklu $i$.

\subsection{Krok 5. Znajdowanie $F_{\mathrm{ref}}(i)$ -- gdzie ,,faza siły''?}

\textbf{Idea:} jeśli $F(t)$ jest cyklem o stałej średniej $F_{st}$, to
w każdym cyklu przecina poziom $F_{st}$ \emph{dwukrotnie}: raz idąc do
góry, raz idąc w dół. Średnia z tych dwóch chwil daje nam ,,środek''
siły w cyklu -- analog \emph{maksimum siły}, ale dokładniej
wyznaczalny numerycznie.

\textbf{Implementacja krok po kroku:}
\begin{enumerate}
  \item Wytnij fragment $F$ w okienku $[ST_i,\,ST_{i+1}]$.
  \item Oblicz amplitudę cyklu $\Delta F = \max F - \min F$.
  \item Sprawdź, czy $F_{st}$ leży w środkowych 50\% zakresu (między
    $F_{st,\,lo} = \min F + 0{,}25\Delta F$ a
    $F_{st,\,hi} = \max F - 0{,}25\Delta F$). Jeśli nie -- cykl
    \emph{odrzucamy} (zbyt zniekształcony).
  \item Znajdź pierwsze przecięcie rosnące $X_{F\uparrow}(i)$ i pierwsze
    przecięcie opadające $X_{F\downarrow}(i)$ z poziomem $F_{st}$
    (interpolacja liniowa).
  \item Wyznacz:
\end{enumerate}
\begin{equation}
  F_{\mathrm{ref}}(i) = \dfrac{1}{2}\bigl(X_{F\uparrow}(i) + X_{F\downarrow}(i)\bigr).
  \label{eq:Fref}
\end{equation}

\textbf{Interpolacja liniowa} przecięcia z poziomem $F_{st}$:
\begin{equation}
  X = t_k + \dfrac{F_{st} - F_k}{F_{k+1} - F_k}(t_{k+1} - t_k).
\end{equation}

\begin{intuicja}
Jeśli sygnał $F$ na próbce $t_k$ jest \emph{poniżej} $F_{st}$, a na
$t_{k+1}$ \emph{powyżej}, to gdzieś po drodze przeszedł przez $F_{st}$.
,,Po drodze'' znaczy: gdy frakcja $(F_{st} - F_k)/(F_{k+1} - F_k)$ pokona
odcinek $[t_k, t_{k+1}]$. To zwykła proporcja.
\end{intuicja}

\subsection{Krok 6. Obliczenie $\varphi(i)$}

Mamy już $TOP_p(i)$ i $F_{\mathrm{ref}}(i)$. Pora policzyć fazę:
\begin{equation}
  \boxed{\varphi(i) = \dfrac{2\pi\bigl(F_{\mathrm{ref}}(i) - TOP_p(i)\bigr)}{\mathrm{Period}(i)}\;\;[\mathrm{rad}].}
  \label{eq:phi_i}
\end{equation}

\begin{intuicja}
\textbf{Co tu się dzieje?} Różnica $F_{\mathrm{ref}}(i) - TOP_p(i)$ to
\textbf{ile sekund} średnia siły wypada \emph{po} maksimum płyty.
Dzieląc przez okres $\mathrm{Period}(i)$, dostajemy ułamek pełnego
cyklu. Mnożąc przez $2\pi$ -- zamieniamy ułamek cyklu na radiany kąta
fazowego. Im większe $\varphi$, tym większe spóźnienie siły -- tym
lepszy amortyzator.
\end{intuicja}

Następnie sprowadzamy wynik do zakresu $[0,\,180^\circ]$:
\begin{equation}
  \varphi_{\deg}(i) = \begin{cases}
    \theta \,\bmod\, 360^\circ, & \text{jeśli wynik } \leq 180^\circ, \\
    360^\circ - (\theta \,\bmod\, 360^\circ), & \text{jeśli wynik } > 180^\circ.
  \end{cases}
\end{equation}

\begin{intuicja}
Faza jest \emph{cykliczna}: $370^\circ$ to to samo co $10^\circ$.
Norma wymaga, by raportować wartości w $[0, 180]$, bo $190^\circ$
to taka sama ,,jakość'' fazy jak $170^\circ$ -- różnica $190 - 180 = 10$
i $180 - 170 = 10$ od antyfazy. Stąd to ,,odbicie'' przy 180.
\end{intuicja}

\subsection{Krok 7. $\varphi_{\min}$ -- jedna liczba na pojazd}

Spośród wszystkich $\varphi(i)$ wybieramy tylko te z cykli o
częstotliwości $f(i) \in [6,\,18]$~Hz. Z nich bierzemy \emph{minimum}:
\begin{equation}
  \varphi_{\min} = \min\{\varphi(i) \colon f(i) \in [\mathrm{MinCalcFreq},\,\mathrm{MaxCalcFreq}]\}.
\end{equation}

Pamiętamy również częstotliwość, przy której to minimum wystąpiło:
\begin{equation}
  f_{\varphi_{\min}} = f(i^*),\quad i^* = \arg\min_i \varphi(i).
\end{equation}

\subsection{Diagram blokowy całej procedury}

\begin{center}
\fbox{\begin{tabular}{c}
  \texttt{[Sygnały t, F, S]} \\
  $\downarrow$ \\
  \texttt{Filtr dolnoprzepustowy 50 Hz (zero-phase)} \\
  $\downarrow$ \\
  \texttt{Detekcja ST\_i (zbocza narastające S)} \\
  $\downarrow$ \\
  \texttt{Macierz kalibracji [delta\_period(i)] (z testu pustej płyty)} \\
  $\downarrow$ \\
  \texttt{Dla każdego cyklu i:} \\
  \texttt{$\ $ -- wyznacz TOP\_p(i) = ST(i) + delta\_period(i)} \\
  \texttt{$\ $ -- znajdź F\_ref(i) (środek dwóch przecięć F z F\_st)} \\
  \texttt{$\ $ -- oblicz phi(i) = 2pi (F\_ref-TOP\_p)/Period} \\
  $\downarrow$ \\
  \texttt{Wybierz min phi(i) dla f(i) w [6, 18 Hz]} \\
  $\downarrow$ \\
  \texttt{Decyzja: phi\_min >= 35 stopni? \texttt{->} sprawny}
\end{tabular}}
\end{center}

% =========================================================================
\section{Model jednomasowy (1 DOF) -- prostszy, do analizy}\label{sec:1dof}
% =========================================================================

\subsection{Po co prostszy model?}

Model 2 DOF jest dokładny, ale ciężko \emph{na nim} analitycznie pokazać,
jak $\varphi_{\min}$ zależy od $c$. Trzeba całkować numerycznie -- co
zajmuje czas i nie daje przejrzystej formuły. Dlatego w trybie ,,Analiza
$\varphi_{\min}(c)$'' używamy uproszczonego modelu \emph{jednomasowego}:

\begin{equation}
  m\,\ddot{x} + c(\dot{x} - \dot{z}) + k(x - z) = 0,\qquad z(t) = -d\cos(\omega t).
  \label{eq:1dof}
\end{equation}

\begin{intuicja}
Bierzemy ,,jedną masę na sprężynie z tłumikiem na wymuszanej
podstawie''. To klasyczny model, który studenci pierwszego roku
mechaniki analizują w pierwszym semestrze. Wymuszenie $z = -d\cos\omega t$
ma minus, bo wygodniej jest tu liczyć przy konwencji, że TOP płyty
występuje przy $\omega t = \pi$.
\end{intuicja}

\subsection{Krok 1. Przeniesienie wymuszenia na prawą stronę}

Rozwijamy nawias w \cref{eq:1dof}:
\begin{equation}
  m\,\ddot{x} - m\ddot{z} + c\dot{x} - c\dot{z} + kx - kz = 0\,??
\end{equation}
Hmm, nie -- wracajmy do oryginału, w \cref{eq:1dof} nie ma $m\ddot{z}$.
Po prostu mnożenie:
\begin{equation}
  m\,\ddot{x} + c\,\dot{x} + k\,x = c\,\dot{z} + k\,z.
\end{equation}

\begin{intuicja}
Lewa strona to ,,reakcja masy ze sprężyną i tłumikiem'' (układ
swobodny). Prawa strona to ,,wymuszenie zewnętrzne'' -- siła, którą
płyta wywiera na układ \emph{przez} sprężynę i tłumik opony.
\end{intuicja}

Podstawiając $z = -d\cos\omega t$, $\dot{z} = d\omega\sin\omega t$:
\begin{equation}
  m\,\ddot{x} + c\,\dot{x} + k\,x = \underbrace{-kd}_{A}\cos\omega t + \underbrace{cd\omega}_{B}\sin\omega t.
\end{equation}

\subsection{Krok 2. Postać próbna -- zgadujemy odpowiedź}

Dla wymuszenia harmonicznego o pulsacji $\omega$ ustalona odpowiedź też
jest harmoniczna o tej samej pulsacji, ale z innym przesunięciem
fazowym:
\begin{equation}
  x(t) = X_c\cos\omega t + X_s\sin\omega t \quad (\text{równoważnie: } X\cos(\omega t - \varphi_x)).
\end{equation}

Wstawmy postać próbną do równania ruchu. Obliczamy pochodne:
\begin{align}
  \dot{x}(t) &= -X_c\omega\sin\omega t + X_s\omega\cos\omega t,\\
  \ddot{x}(t) &= -X_c\omega^2\cos\omega t - X_s\omega^2\sin\omega t.
\end{align}

Wszystko zbieramy ,,po $\cos$'' i ,,po $\sin$'':
\begin{align}
  \text{przy } \cos\omega t:\quad & (k - m\omega^2)X_c + c\omega X_s = -kd, \\
  \text{przy } \sin\omega t:\quad & -c\omega X_c + (k - m\omega^2)X_s = cd\omega.
\end{align}

\begin{intuicja}
Skoro lewa strona pierwotnego równania ma dawać tę samą funkcję czasu
co prawa, to \emph{osobno} musimy dopasować amplitudę przy $\cos\omega t$
oraz amplitudę przy $\sin\omega t$. Dlatego dostajemy dwa równania
liniowe na dwie niewiadome $X_c, X_s$.
\end{intuicja}

\subsection{Krok 3. Rozwiązanie układu liniowego}

Niech $\Delta = (k - m\omega^2)^2 + (c\omega)^2$. Stosując wzory Cramera:
\begin{align}
  X_c &= \dfrac{(k - m\omega^2)(-kd) - (c\omega)(cd\omega)}{\Delta}
       = \dfrac{-d\bigl[k(k - m\omega^2) + (c\omega)^2\bigr]}{\Delta},\\
  X_s &= \dfrac{(k - m\omega^2)(cd\omega) - (c\omega)(-kd)}{\Delta}
       = \dfrac{cd\omega(k - m\omega^2) + cd\omega k}{\Delta}
       = \dfrac{-cd m\omega^3}{\Delta}.
\end{align}

\textbf{Sprawdzenie sensowności:}
\begin{itemize}
  \item Dla $c=0$ (brak tłumienia): $X_s = 0$, więc $x$ jest tylko
    cosinusem -- czyli w fazie z $z$ pod albo nad rezonansem (180° w
    rezonansie -- klasyczny obraz).
  \item Dla bardzo dużego $c$: $\Delta \to (c\omega)^2$, więc
    $X_c \to -d$, $X_s \to -dm\omega/c \to 0$. Czyli $x \to -d\cos\omega t = z$
    -- masa porusza się razem z płytą. To poprawnie odpowiada granicy
    ,,nieskończenie sztywnego sprzężenia''.
\end{itemize}

\subsection{Krok 4. Siła kontaktu i jej faza}

Siła kontaktu opony to:
\begin{equation}
  F_{\mathrm{dyn}}(t) = k(x - z) + c(\dot{x} - \dot{z}).
\end{equation}

Wszystkie składniki ($x, z, \dot{x}, \dot{z}$) to kombinacje $\cos\omega t$
i $\sin\omega t$. Wynikiem jest też kombinacja $A_F\cos\omega t + B_F\sin\omega t$,
czyli harmoniczna o pewnej amplitudzie $|F| = \sqrt{A_F^2 + B_F^2}$ i
fazie $\varphi_F = \arctantwo(B_F, A_F)$.

\subsection{Krok 5. Przesunięcie fazowe $\varphi$}

Faza $z(t) = -d\cos\omega t = d\cos(\omega t - \pi)$ to po prostu $\pi$
(czyli $180^\circ$). Faza $F_{\mathrm{dyn}}$ wynosi $\varphi_F$
(z arctan2). Przesunięcie:
\begin{equation}
  \varphi(\omega) = \varphi_F(\omega) - \pi,
\end{equation}
sprowadzone do zakresu $[0,\,180^\circ]$.

Mając krzywą $\varphi(\omega)$, wystarczy obliczyć jej minimum w
$\omega \in [12\pi,\,36\pi]$ rad/s (czyli $f \in [6,18]$~Hz):
\begin{equation}
  \varphi_{\min}(c) = \min_{f\in[6,\,18]\,\mathrm{Hz}} \varphi(2\pi f;\,c).
\end{equation}

\subsection{Zachowanie graniczne (sanity check)}

\begin{uwaga}[Granica $c\to 0$]
Bez tłumienia układ \emph{nie rozprasza energii}. Faza $F$ względem $z$
to albo $0^\circ$ (pod rezonansem), albo $180^\circ$ (nad). Po
sprowadzeniu do $[0,180]$ minimum to $0^\circ$. Numerycznie potwierdzamy:
$\varphi_{\min}(c=0) \approx 0^\circ$.
\end{uwaga}

\begin{uwaga}[Granica $c\to\infty$]
Tłumik dominuje, układ porusza się sztywno z płytą, faza $F$ jest
$90^\circ$ (charakterystyka dla siły lepkiej $c\dot{z}$). Po
sprowadzeniu minimum jest niezerowe i \emph{rośnie monotonicznie} z $c$.
\end{uwaga}

\begin{przyklad}
Dla $m = 382$~kg, $k = 253\,161$~N/m, $d = 3$~mm:
\begin{center}
\begin{tabular}{rcc}
\toprule
$c$ [N$\cdot$s/m] & $\varphi_{\min}$ [$^\circ$] & $f_{\varphi_{\min}}$ [Hz] \\
\midrule
400  &  6{,}05 & 7{,}11 \\
800  & 12{,}05 & 7{,}11 \\
1200 & 17{,}97 & 7{,}14 \\
1474 & 21{,}95 & 7{,}17 \\
2000 & 29{,}38 & 7{,}23 \\
3000 & 42{,}54 & 7{,}44 \\
\bottomrule
\end{tabular}
\end{center}
$\varphi_{\min}$ rośnie monotonicznie z $c$, a próg $35^\circ$ jest
przekraczany pomiędzy $c=2000$ a $c=3000$.
\end{przyklad}

% =========================================================================
\section{Wskaźniki diagnostyczne}\label{sec:wskazniki}
% =========================================================================

\subsection{Statyczne obciążenie koła $F_{st}$}

\begin{equation}
  F_{st} = (M + m)\,g.
\end{equation}

\begin{intuicja}
To po prostu \emph{ciężar} jednej ćwiartki auta -- masa razy
przyspieszenie ziemskie. Mierzony przed testem (przy nieruchomej
płycie). Wzorcowy punkt odniesienia dla większości obliczeń.
\end{intuicja}

\subsection{Ekstrema siły w cyklu}

\begin{align}
  \mathrm{maxF}(i) &= \max_{t\in\mathrm{Period}(i)} F(t),\\
  \mathrm{minF}(i) &= \min_{t\in\mathrm{Period}(i)} F(t),\\
  \Delta F(i) &= \mathrm{maxF}(i) - \mathrm{minF}(i).
\end{align}

I globalne wartości w paśmie 6--18~Hz:
\begin{equation}
  F_{\min} = \min_{i} \mathrm{minF}(i), \quad F_{\max} = \max_{i} \mathrm{maxF}(i).
\end{equation}

\subsection{Maksymalna amplituda dynamiczna $FA_{\max}$}

\begin{equation}
  FA_{\max} = \begin{cases}
    F_{st} - F_{\min}, & F_{\min} \geq F_{\mathrm{UnderLim}} \\
    F_{\max} - F_{st}, & F_{\min} < F_{\mathrm{UnderLim}}, F_{\max} \leq F_{\mathrm{OverLim}} \\
    \max(F_{\mathrm{OverLim}} - F_{st},\,F_{st} - F_{\mathrm{UnderLim}}), & \text{oba przypadki}
  \end{cases}
\end{equation}

\begin{intuicja}
Trzy przypadki obsługują różne sytuacje:
\begin{enumerate}
  \item \emph{Wszystko OK} (siła nie spada do zera): patrzymy, jak głęboko
    siła zeszła poniżej ciężaru statycznego.
  \item \emph{Siła chwilowo zerowa} (koło odrywa się): tej części nie
    da się dokładnie zmierzyć, więc patrzymy na górną odbitkę.
  \item \emph{Skrajne przypadki saturacji}: bierzemy maksimum z obu.
\end{enumerate}
\end{intuicja}

\subsection{Względna amplituda $RFA_{\max}$}

\begin{equation}
  RFA_{\max} = \dfrac{FA_{\max}}{F_{st}}\cdot 100\,\%.
  \label{eq:rfa}
\end{equation}

To samo co $FA_{\max}$, ale w procentach -- dzięki temu wartość nie
zależy od ciężaru auta i można porównywać różne pojazdy.

\subsection{Częstotliwość rezonansowa $f_{\mathrm{res}}$}

\begin{equation}
  f_{\mathrm{res}} = f(i^*), \quad i^* = \arg\max_i FA(i).
\end{equation}

To częstotliwość okresu, w którym wystąpiła największa zmiana siły.
Powinna być bliska teoretycznej $f_2 \approx 13{,}5$~Hz.

\subsection{Amplituda statyczna $H_{25}$ -- pomiar opony}

W fazie stabilizacji ($f = 25$~Hz) uśredniamy amplitudę peak-to-peak
po około 10 okresach:
\begin{equation}
  H_{25} = \dfrac{1}{2}\bigl\langle \Delta F(i)\bigr\rangle_{f(i)\approx 25\,\mathrm{Hz}}.
\end{equation}

\begin{intuicja}
Przy 25~Hz układ jest powyżej obu rezonansów -- praktycznie wszystko
\emph{nieruchome} oprócz opony, która przenosi drgania. Stąd $H_{25}$
mówi nam właściwie o sztywności opony, a nie amortyzatora.
\end{intuicja}

\subsection{Sztywność opony $\mathrm{rig}$}

\begin{equation}
  \mathrm{rig} = a_{\mathrm{rig}}\,\dfrac{H_{25}}{e_p} + b_{\mathrm{rig}},\qquad
  a_{\mathrm{rig}} = 0{,}571,\; b_{\mathrm{rig}} = 46{,}0,
\end{equation}
gdzie $e_p$ w milimetrach.

Granice ostrzeżeń (rozdz.~5.5 normy):
\begin{itemize}
  \item $\mathrm{rig} < 160$~N/mm $\Rightarrow$ opona niedopompowana,
  \item $\mathrm{rig} > 400$~N/mm $\Rightarrow$ opona przepompowana.
\end{itemize}

\subsection{Wskaźnik EUSAMA}

\begin{equation}
  \mathrm{EUSAMA} = \dfrac{F_{\min}}{F_{st}}\cdot 100\,\% \quad\text{w paśmie 6--18 Hz}.
\end{equation}

\textbf{Status w normie SPECSUS2018:} \emph{tylko informacyjny}. W naszej
aplikacji jednak \emph{wykorzystujemy go praktycznie} jako wskaźnik
pomocniczy -- patrz \cref{sec:diag_app}.

% =========================================================================
\section{Kryteria pass/fail i logika diagnostyczna aplikacji}\label{sec:kryteria}\label{sec:diag_app}
% =========================================================================

\subsection{Kryterium absolutne normy}

Sama norma definiuje próg:
\begin{equation}
  \mathrm{AC}_{\varphi_{\min}} = 35^\circ.
\end{equation}

I trzy stany:
\begin{itemize}
  \item $\varphi_{\min} \geq 35^\circ \Rightarrow$ \textbf{sprawny};
  \item $25^\circ \leq \varphi_{\min} < 35^\circ \Rightarrow$ stan graniczny;
  \item $\varphi_{\min} < 25^\circ \Rightarrow$ \textbf{niesprawny}.
\end{itemize}

\subsection{Dlaczego potrzebujemy wskaźników dodatkowych?}

\begin{uwagabox}
W modelu ćwiartkowym (2 DOF) okazuje się, że $\varphi_{\min}$ \emph{nie
zawsze spada poniżej $35^\circ$} dla uszkodzonego amortyzatora, jeśli
tłumienie opony $c_m$ jest nominalne ($\sim 150$ N$\cdot$s/m). Rezonans
masy nieresorowanej zostaje wówczas wytłumiony przez samą oponę,
a wykrywany kąt $\varphi(f)$ w pobliżu rezonansu pozostaje w okolicach
$80$--$90^\circ$ niezależnie od $c_M$.

Sytuacja jest natomiast jednoznaczna w sygnale samej siły: gdy $c_M$
spada, układ wpada w silny rezonans, $F(t)$ spada \textbf{poniżej zera}
(koło odrywa się od płyty) i \emph{EUSAMA} spada do wartości ujemnych.
\end{uwagabox}

Stąd nasza wielowarstwowa logika decyzyjna:

\begin{center}
\fbox{\begin{tabular}{l}
  \texttt{IF F\_min < 0 (koło odrywa się)} \\
  \texttt{$\quad$ -> NIESPRAWNY} \\
  \texttt{ELSE IF EUSAMA < 30\%} \\
  \texttt{$\quad$ -> NIESPRAWNY} \\
  \texttt{ELSE IF phi\_min < 25} \\
  \texttt{$\quad$ -> NIESPRAWNY} \\
  \texttt{ELSE IF phi\_min < 35} \\
  \texttt{$\quad$ -> STAN GRANICZNY} \\
  \texttt{ELSE IF EUSAMA < 45\%} \\
  \texttt{$\quad$ -> STAN GRANICZNY} \\
  \texttt{ELSE -> SPRAWNY}
\end{tabular}}
\end{center}

\subsection{Symetria osi (kryterium względne)}

Aplikacja pozwala porównać oba koła tej samej osi:
\begin{equation}
  D\varphi_{\min} = \dfrac{|\varphi_{\min,L} - \varphi_{\min,R}|}{\max(\varphi_{\min,L},\,\varphi_{\min,R})}\cdot 100\,\%.
\end{equation}
Norma dopuszcza $RC_{\varphi_{\min}} = 30\,\%$ (rozdz.~5.6).

\subsection{Tolerancje powtarzalności (rozdz.~6.1.4 normy)}

\begin{center}
\begin{tabular}{lccc}
\toprule
\textbf{Wielkość} & \textbf{Powtarzalność} & \textbf{Powtarzalność (pozycje)} & \textbf{Błąd całkowity} \\
\midrule
$\varphi_{\min} > 30^\circ$ & $\pm 3^\circ$  & $\pm 4{,}5^\circ$ & $7{,}5^\circ$ \\
$\varphi_{\min} = 0^\circ$  & $\pm 6^\circ$  & $\pm 9^\circ$     & $15^\circ$ \\
$RFA_{\max}$                & $\pm 1{,}5\,\%$ & $\pm 3\,\%$      & $5\,\%$ \\
$H_{25}$                    & $\pm 24$~daN    & ---              & $8\,\%$ \\
$F_{st} \geq 300$~daN       & $\pm 2\,\%$     & ---              & $\pm 2\,\%$ \\
$F_{st} < 300$~daN          & $\pm 6$~daN     & ---              & $\pm 6$~daN \\
\bottomrule
\end{tabular}
\end{center}

% =========================================================================
\section{Implementacja numeryczna w Pythonie}\label{sec:kod}
% =========================================================================

\subsection{Generowanie wymuszenia $z(t)$ (model 2 DOF)}

\begin{lstlisting}[style=py, caption={Cykliczne generowanie $z(t)$, $\dot{z}(t)$, $f(t)$.}]
def _generuj_wymuszenie(self, t):
    z = np.zeros_like(t); z_dot = np.zeros_like(t)
    f_step = np.zeros_like(t); st_indices = []
    i, t_cycle_start = 0, t[0]
    f_cycle = self.czestotliwosc_docelowa(t_cycle_start)
    while i < len(t):
        period = 1.0 / f_cycle
        t_cycle_end = t_cycle_start + period
        if not st_indices or st_indices[-1] != i:
            st_indices.append(i)
        while i < len(t) and t[i] < t_cycle_end:
            theta = 2*np.pi*(t[i] - t_cycle_start)/period
            theta_dot = 2*np.pi/period
            z[i] = self.d * np.cos(theta)
            z_dot[i] = -self.d * np.sin(theta) * theta_dot
            f_step[i] = f_cycle
            i += 1
        if i >= len(t): break
        t_cycle_start = t[i]
        f_cycle = self.czestotliwosc_docelowa(t_cycle_start)
    return z, z_dot, f_step, st_indices
\end{lstlisting}

\textbf{Co tu się dzieje?} Pętla zewnętrzna iteruje po cyklach (każdy
o długości $1/f$). W obrębie cyklu fazę $\theta$ liczymy
proporcjonalnie do tego, jak daleko ,,jesteśmy'' od startu cyklu. Nie
liczymy $\theta = 2\pi f \cdot t$ globalnie, bo $f$ się zmienia -- to
prowadziłoby do nieciągłości.

\subsection{Całkowanie równań ruchu (ODE)}

Stosujemy \texttt{scipy.integrate.odeint} (algorytm LSODA), który
automatycznie przełącza się między metodą Adamsa (układ ,,łagodny'')
a BDF (układ ,,sztywny''). Krok adaptacyjny zapewnia stabilność nawet
przy ostrych rezonansach.

\begin{lstlisting}[style=py, caption={Funkcja prawej strony układu \cref{eq:stan}.}]
def f_rhs(X, t_val):
    idx = int(round(t_val * fs))
    idx = max(0, min(idx, len(z) - 1))
    B = np.array([0.0, 0.0,
                  (self.km*z[idx] + self.cm*z_dot[idx])/self.m,
                  0.0])
    return self.A @ X + B
\end{lstlisting}

\subsection{Obliczanie $\varphi(i)$}

\begin{lstlisting}[style=py, caption={Wyznaczenie $\varphi(i)$ z procedury MinPhaseShift.}]
for i, ((a, b), v) in enumerate(zip(pairs, valid)):
    if not v: continue
    period = t[b] - t[a]
    dt_phase = F_ref_t[i] - TOPp[i]
    phi_rad = 2 * np.pi * dt_phase / period
    phi_deg_val = np.degrees(phi_rad) % 360.0
    if phi_deg_val > 180:
        phi_deg_val = 360 - phi_deg_val
    phi_deg[i] = phi_deg_val
\end{lstlisting}

\subsection{Detekcja przecięć $F = F_{st}$}

\begin{lstlisting}[style=py, caption={Detekcja przecięć $F(t) = F_{st}$.}]
signs = np.sign(seg_F - F_st)
crossings = np.where(np.diff(signs) != 0)[0]
for k in crossings:
    f1, f2 = seg_F[k], seg_F[k+1]
    frac = (F_st - f1) / (f2 - f1)
    t_cross = seg_t[k] + frac * (seg_t[k+1] - seg_t[k])
    if f1 < F_st <= f2 and x_up is None:
        x_up = t_cross
    elif f1 >= F_st > f2 and x_dn is None:
        x_dn = t_cross
F_ref_t[i] = 0.5 * (x_up + x_dn)
\end{lstlisting}

\textbf{Wyjaśnienie:} \texttt{np.sign} zamienia wartości na $-1, 0, +1$
zależnie od znaku $F - F_{st}$. \texttt{np.diff} liczy różnice
między kolejnymi próbkami -- niezerowa różnica oznacza, że znak się
zmienił, czyli funkcja przecięła zero (poziom $F_{st}$). Dla każdego
przecięcia interpolujemy liniowo do dokładnego czasu i klasyfikujemy
jako rosnące/opadające na podstawie znaków przed i po.

% =========================================================================
\section{Wszystkie parametry stałe normy (rozdz.~9)}
% =========================================================================

\begin{center}
\begin{tabular}{lll}
\toprule
\textbf{Symbol} & \textbf{Wartość} & \textbf{Opis} \\
\midrule
$\mathrm{AC}_{\varphi_{\min}}$ & $35^\circ$ & Absolutne kryterium $\varphi_{\min}$ \\
$\mathrm{MaxCalcFreq}$  & 18~Hz   & Górna granica zakresu pomiarowego \\
$\mathrm{MinCalcFreq}$  & 6~Hz    & Dolna granica zakresu pomiarowego \\
$\Delta T_{\mathrm{meas}}$ & 7{,}5~s & Minimalny czas rampy 18$\to$6 Hz \\
$\Delta TfMaxSlope$     & 3~Hz/s  & Maks.~nachylenie zjazdu $f$ \\
$\Delta TfLinErr$       & 2~Hz    & Dopuszczalny błąd liniowości $f$ \\
$h_{PS}$                & 6~mm    & Amplituda peak-to-peak płyty \\
$\mathrm{RFstFMax}$     & 25~\%   & Margines górny detekcji $F_{\mathrm{ref}}$ \\
$\mathrm{RFstFMin}$     & 25~\%   & Margines dolny detekcji $F_{\mathrm{ref}}$ \\
$\mathrm{FUnderLimPerc}$ & 1~\%   & Próg detekcji underflow $F(t)$ \\
$\mathrm{DeltaF}$        & 5~Hz   & Margines detekcji $\varphi_{\min}$ pod $f_{\mathrm{res}}$ \\
$a_{\mathrm{rig}}$       & 0{,}571 & Współ.~liniowy wzoru sztywności opony \\
$b_{\mathrm{rig}}$       & 46{,}0  & Wyraz wolny wzoru sztywności opony \\
$\mathrm{rigLoLim}$      & 160~N/mm & Dolna granica ostrzeżenia o ciśnieniu \\
$\mathrm{rigHiLim}$      & 400~N/mm & Górna granica ostrzeżenia o ciśnieniu \\
$\mathrm{StatWLim}$      & 25~daN  & Maks.~zmiana $F_{st}$ przed/po teście \\
$F_{\mathrm{UnderLim}}$  & $F_{st}\cdot 1\,\%$ & Próg underflow sygnału \\
\bottomrule
\end{tabular}
\end{center}

% =========================================================================
\section{Domyślne parametry pojazdu (oś przednia M1)}
% =========================================================================

\begin{center}
\begin{tabular}{lll}
\toprule
\textbf{Parametr} & \textbf{Wartość} & \textbf{Jednostka} \\
\midrule
$M$ -- masa resorowana   & 346     & kg \\
$m$ -- masa nieresorowana & 36      & kg \\
$k_M$ -- sztywność zawieszenia & 25\,570 & N/m \\
$k_m$ -- sztywność opony      & 253\,161 & N/m \\
$c_M$ -- tłumienie amortyzatora (sprawny) & 1\,474 & N$\cdot$s/m \\
$c_m$ -- tłumienie opony     & 150     & N$\cdot$s/m \\
$d$  -- amplituda płyty       & 0{,}003 & m \\
$m_p$ -- masa pustej płyty    & 20      & kg \\
$F_{st}$ -- statyczne obciążenie & 3\,747 & N \\
\bottomrule
\end{tabular}
\end{center}

% =========================================================================
\section{Walidacja i wyniki referencyjne}
% =========================================================================

\subsection{Test symulacji domyślnej (zdrowy amortyzator)}

Dla parametrów domyślnych, symulacja z $f_s = 2$~kHz daje:
\begin{align}
  \varphi_{\min}   &\approx 88{,}9^\circ \quad @ \quad f \approx 12{,}0~\mathrm{Hz}, \\
  \mathrm{EUSAMA}  &\approx 57{,}0\,\%, \\
  RFA_{\max}       &\approx 43{,}1\,\%, \\
  H_{25}           &\approx 1030~\mathrm{N},\\
  \mathrm{rig}     &\approx 242~\mathrm{N/mm}.
\end{align}

\textbf{Wniosek}: amortyzator z $c_M = 1474$~N$\cdot$s/m jest
\emph{sprawny} ($\varphi_{\min} \gg 35^\circ$ \emph{oraz} EUSAMA $> 45$\%);
opona prawidłowo dopompowana ($\mathrm{rig}\in[160,\,400]$).

\subsection{Test symulacji z uszkodzonym amortyzatorem}

Dla $c_M = 100$, $c_m = 30$:
\begin{align}
  \varphi_{\min} &\approx 43^\circ, \\
  \mathrm{EUSAMA} &\approx -228\,\% \quad (!) \\
  F_{\min} &\approx -8540~\mathrm{N} \quad (\text{koło odrywa się}).
\end{align}

\textbf{Wniosek aplikacji}: NIESPRAWNY (koło traci kontakt) -- pomimo
że samo $\varphi_{\min} = 43^\circ > 35^\circ$ sugerowałoby ,,sprawny''.
To dokładnie ta sytuacja, dla której wprowadzono wieloskładnikowe
kryterium z \cref{sec:diag_app}.

% =========================================================================
\section{Architektura kodu źródłowego}
% =========================================================================

\subsection{Pliki źródłowe}

\begin{description}
  \item[\texttt{egea\_suspension\_model.py}] -- klasy
    \texttt{ModelCwiartkowy} (2 DOF) i \texttt{ModelJednomasowy} (1 DOF)
    oraz parametry normy w \texttt{ParametryEGEA}.
  \item[\texttt{phase\_shift.py}] -- algorytm wyznaczania $\varphi(i)$,
    detekcja $ST_i$, kalibracja $\Delta\mathrm{Period}(i)$, obliczanie
    $\varphi_{\min}$, $RFA_{\max}$, $H_{25}$, $\mathrm{rig}$.
  \item[\texttt{app.py}] -- aplikacja Streamlit (interfejs) w trzech
    trybach:
    \begin{enumerate}
      \item symulacja 2 DOF -- generowanie syntetycznego pomiaru
        i diagnostyka;
      \item import CSV -- analiza rzeczywistego pomiaru;
      \item analiza $\varphi_{\min}(c)$ -- model 1 DOF, krzywa
        parametryczna.
    \end{enumerate}
\end{description}

\subsection{Zależności (wymagane biblioteki)}

\begin{itemize}
  \item \texttt{numpy} -- algebra liniowa i przetwarzanie tablic,
  \item \texttt{scipy} -- całkowanie ODE (\texttt{odeint}, LSODA),
    filtry (\texttt{butter}, \texttt{filtfilt}),
  \item \texttt{pandas} -- import/eksport CSV,
  \item \texttt{plotly} -- interaktywne wykresy,
  \item \texttt{streamlit} -- interfejs WWW.
\end{itemize}

\subsection{Jak uruchomić}

\begin{lstlisting}[style=py, language=bash, caption={Skrypty uruchomieniowe.}]
# Instalacja
pip install -r requirements.txt

# Aplikacja webowa
streamlit run app.py

# Test modułów (bez UI)
python3 egea_suspension_model.py
python3 phase_shift.py
\end{lstlisting}

% =========================================================================
\section{Format pliku CSV}
% =========================================================================

Plik wejściowy w trybie ,,Import pomiaru'' składa się z trzech kolumn:
\begin{center}
\begin{tabular}{lll}
\toprule
\textbf{Kolumna} & \textbf{Typ} & \textbf{Opis} \\
\midrule
$t$ & float [s] & czas próbki, równomiernie spróbkowany \\
$F$ & float [N] & siła kontaktu opony $F(t)$ (po odjęciu $F_p$) \\
$S$ & int 0/1   & sygnał czujnika płyty (1 w $ST_i$) \\
\bottomrule
\end{tabular}
\end{center}

\textbf{Wymagania techniczne:}
\begin{itemize}
  \item Częstotliwość próbkowania $f_s \geq 2 f_{\max} = 36$~Hz
    (twierdzenie Nyquista); zalecane $f_s \geq 500$~Hz dla wiarygodnej
    detekcji przecięć $F_{\mathrm{ref}}$.
  \item $F_{st}$ wyznaczone przed testem lub automatycznie z
    pierwszych $0{,}5$~s.
  \item Separator: przecinek lub średnik (auto-detekcja).
\end{itemize}

% =========================================================================
\section{Dyskusja i ograniczenia}
% =========================================================================

\subsection{Założenia upraszczające w modelu}

\begin{itemize}
  \item \textbf{Liniowość:} zakładamy, że sprężyny i tłumiki mają
    charakterystyki liniowe (siła proporcjonalna do położenia/prędkości).
    Rzeczywiste elementy mają nieliniowości: suche tarcie, ograniczniki
    skoku, progresywność.
  \item \textbf{Brak interakcji osi:} modelujemy jedno koło. W rzeczywistym
    pojeździe drgania jednej osi częściowo przenoszą się na drugą przez
    nadwozie.
  \item \textbf{Sztywność płyty:} traktujemy płytę jako idealnie
    sztywną. W praktyce ma własną podatność.
  \item \textbf{Brak luzów:} pomijamy luzy w łączach mechanicznych.
\end{itemize}

Rozbieżność między $\varphi_{\min}^{\text{model}}$ a
$\varphi_{\min}^{\text{rzeczywisty}}$ wynosi typowo kilka stopni.

\subsection{Czułość $\varphi_{\min}$ na tłumienie}

W modelu 2 DOF $\varphi_{\min}$ silnie zależy od \emph{sumarycznego}
tłumienia w pobliżu rezonansu masy nieresorowanej. Dominującą rolę
odgrywa $c_m$ (tłumienie opony) -- to dlatego nasza aplikacja korzysta
też z EUSAMA i wykrywania odrywania koła.

W modelu 1 DOF $\varphi_{\min}$ wzrasta wprost z $c$ -- to dlatego
w analizie parametrycznej używamy 1 DOF dla klarowności obrazu.

\subsection{Konieczność kalibracji}

Bez kalibracji dynamicznej $\Delta\mathrm{Period}(i)$ nie jest znane,
co prowadzi do błędu fazy. W aplikacji domyślnie przyjmujemy
$\Delta\mathrm{Period} = 0$ (założenie sprzętowe, że $ST_i \approx TOP$).
Dla profesjonalnych pomiarów wymagana jest macierz kalibracyjna.

% =========================================================================
\begin{thebibliography}{9}

\bibitem{EGEA2019}
EGEA SWG6, \emph{Suspension tester specifications SPECSUS2018},
European Garage Equipment Association, luty 2019.

\bibitem{Tsymberov1996}
A.~Tsymberov, \emph{An improved Non-Intrusive Automotive Suspension Testing
Apparatus with Means to Determine the Condition of a Dampers},
SAE Technical Paper 960735, Hunter Engineering Co., 1996.

\bibitem{Klapka2017}
M.~Klapka, I.~Mazurek, O.~Macháček, M.~Kubík,
\emph{Twilight of the EUSAMA diagnostic methodology},
Meccanica 52(9), 2017, s.~2023--2034.\\
DOI: 10.1007/s11012-016-0566-0.

\bibitem{EUSAMA1976}
EUSAMA, \emph{Recommendation for performance test specification of on-car
vehicle suspension testing system}, EUSAMA-TS-02-76, 1976.

\bibitem{Calvo2005}
J.A.~Calvo, V.~Diaz, J.L.~San Roman,
\emph{Establishing inspection criteria to verify the dynamic behaviour
of the vehicle suspension system by a platform vibrating test bench},
Int.\ J.\ Vehicle Design, 38(4), 2005.

\bibitem{Kaiser1977}
J.F.~Kaiser, W.A.~Reed, \emph{Data smoothing using low-pass filters},
Review of Scientific Instruments, 48(11), 1977, s.~1447--1457.

\bibitem{Balsarotti2000}
S.~Balsarotti, W.~Bradley,
\emph{Experimental Evaluation of a Non-Intrusive Automotive Suspension
Testing Apparatus}, SAE Technical Paper SAE 2000-01-1329,
Hunter Engineering Co.\ \& General Motors Corp., 2000.

\end{thebibliography}

\appendix
\clearpage

% =========================================================================
\section{Dodatek A: Pełna lista symboli}
% =========================================================================

\begin{longtable}{lll}
\toprule
\textbf{Symbol} & \textbf{Opis} & \textbf{Jednostka} \\
\midrule
\endhead
$M$        & Masa resorowana                    & kg \\
$m$        & Masa nieresorowana                 & kg \\
$m_p$      & Masa pustej płyty stanowiska       & kg \\
$k_M$      & Sztywność zawieszenia              & N/m \\
$k_m$      & Sztywność opony                    & N/m \\
$c_M$      & Tłumienie amortyzatora             & N$\cdot$s/m \\
$c_m$      & Tłumienie opony                    & N$\cdot$s/m \\
$d$        & Amplituda płyty                    & m \\
$e_p$      & $d$ w mm (do wzoru rig)            & mm \\
$h_{PS}$   & Peak-to-peak płyty $= 2d$          & mm \\
$g$        & Przyspieszenie ziemskie 9{,}81     & m/s$^2$ \\
$f(t)$     & Docelowa częstotliwość             & Hz \\
$f(i)$     & Częstotliwość okresu $i$           & Hz \\
$f_{\mathrm{res}}$ & Częstotliwość rezonansowa  & Hz \\
$f_{\varphi_{\min}}$ & Częstotliwość minimum $\varphi$ & Hz \\
$\omega$   & Pulsacja, $\omega = 2\pi f$        & rad/s \\
$\theta(t)$ & Faza wymuszenia                   & rad \\
$z(t)$     & Położenie płyty                    & m \\
$\dot{z}(t)$ & Prędkość płyty                   & m/s \\
$\ddot{z}(t)$ & Przyspieszenie płyty            & m/s$^2$ \\
$x_m(t)$   & Położenie masy nieresorowanej      & m \\
$x_M(t)$   & Położenie masy resorowanej         & m \\
$F(t)$     & Siła kontaktu opony                & N \\
$F_p(t)$   & Siła pustej płyty                  & N \\
$F_r(t)$   & Surowy sygnał siły z tensometrów   & N \\
$F_{st}$   & Statyczne obciążenie koła          & N \\
$F_{\min},F_{\max}$ & Globalne ekstrema $F(t)$  & N \\
$F_{\mathrm{ref}}(i)$ & Punkt odniesienia okresu & s \\
$ST(i)$    & Znacznik wyzwalania                & s \\
$TOP_p(i)$ & TOP płyty (skorygowane)            & s \\
$\Delta\mathrm{Period}(i)$ & Korekcja kalibracyjna & s \\
$\mathrm{Period}(i)$ & Długość okresu $i$       & s \\
$FA_{\max}$ & Maks.~amplituda dynamiczna        & N \\
$RFA_{\max}$ & Względna amplituda maks.         & \% \\
$H_{25}$   & Amplituda statyczna przy 25 Hz     & N \\
$\mathrm{rig}$ & Sztywność opony                & N/mm \\
$\varphi(i)$ & Przesunięcie fazowe okresu $i$   & ${}^\circ$ \\
$\varphi_{\min}$ & Minimum $\varphi$ w 6--18 Hz & ${}^\circ$ \\
$\varphi_{\max}$ & $\varphi$ przy 18 Hz         & ${}^\circ$ \\
$\mathrm{AC}_{\varphi_{\min}}$ & Kryterium absolutne & ${}^\circ$ \\
$\mathrm{EUSAMA}$ & Wskaźnik historyczny        & \% \\
$\Delta T_{25}$ & Czas stabilizacji 25 Hz       & s \\
$\Delta T_{\mathrm{meas}}$ & Czas rampy pomiarowej & s \\
$\mathrm{MaxCalcFreq}$ & Górna granica $f$ pomiaru & Hz \\
$\mathrm{MinCalcFreq}$ & Dolna granica $f$ pomiaru & Hz \\
$\mathrm{RFstFMax/Min}$ & Marginesy detekcji $F_{\mathrm{ref}}$ & \% \\
\bottomrule
\end{longtable}

% =========================================================================
\section{Dodatek B: Przykładowy raport tekstowy z aplikacji}
% =========================================================================

\begin{lstlisting}[basicstyle=\ttfamily\footnotesize, frame=single,
                   backgroundcolor=\color{codebg}, breaklines=true]
RAPORT BADANIA STANOWISKA EGEA
================================

PARAMETRY POJAZDU
  M  = 346.0 kg
  m  = 36.0 kg
  kM = 25570.0 N/m
  km = 253161.0 N/m
  cM = 1474.0 N*s/m
  cm = 150.0 N*s/m
  d  = 3.00 mm

WYNIKI
  F_st     = 3747.4 N
  phi_min  = 88.86 stopni
  f_phi_min= 11.98 Hz
  EUSAMA   = 57.00 %
  RFA_max  = 43.06 %
  H25      = 1029.7 N
  rig      = 242.0 N/mm

OCENA
  Kryterium AC_phi_min = 35 stopni
  Stan amortyzatora: SPRAWNY
\end{lstlisting}

% =========================================================================
\section{Dodatek C: Pełne wyprowadzenie odpowiedzi 1 DOF (krok po kroku)}
% =========================================================================

Tu rozwijamy pełne wyprowadzenie z \cref{sec:1dof} w drobiazgowy sposób.
Każdy krok jest opatrzony notatką pedagogiczną.

\subsection*{Krok C.1: Wyjściowe równanie}

\begin{equation}
  m\,\ddot{x} + c(\dot{x} - \dot{z}) + k(x - z) = 0,\qquad z(t) = -d\cos(\omega t).
  \tag{1DOF}
\end{equation}

\subsection*{Krok C.2: Postać próbna $x = X_c\cos + X_s\sin$}

\begin{equation}
  x(t) = X_c\cos\omega t + X_s\sin\omega t.
\end{equation}

\subsection*{Krok C.3: Pochodne}

\begin{align}
  \dot{x}(t) &= -X_c\omega\sin\omega t + X_s\omega\cos\omega t,\\
  \ddot{x}(t) &= -X_c\omega^2\cos\omega t - X_s\omega^2\sin\omega t.
\end{align}

\subsection*{Krok C.4: Podstawiamy i porównujemy współczynniki}

Po podstawieniu i wyrugowaniu $\ddot z$ (równanie pierwotne nie zawiera
$\ddot z$, tylko $\dot z$ i $z$), zbieramy współczynniki przy $\cos\omega t$:
\begin{equation}
  -mX_c\omega^2 + c\omega X_s + kX_c - k(-d) = 0,
\end{equation}
co po porządkowaniu daje:
\begin{equation}
  (k - m\omega^2)X_c + c\omega X_s = -kd. \tag{C-1}
\end{equation}

Analogicznie współczynniki przy $\sin\omega t$:
\begin{equation}
  -mX_s\omega^2 - c\omega X_c + kX_s - c\omega(-d\omega)/(\omega) = 0,
\end{equation}
hmm -- najprościej zauważyć od razu: $-c\dot z = -c(d\omega\sin\omega t)$
ma znak \emph{ujemny}, a po przeniesieniu na prawą daje $+cd\omega\sin\omega t$.
Po uporządkowaniu:
\begin{equation}
  -c\omega X_c + (k - m\omega^2)X_s = cd\omega. \tag{C-2}
\end{equation}

\subsection*{Krok C.5: Wyznacznik układu}

\begin{equation}
  \Delta = (k - m\omega^2)^2 + (c\omega)^2.
\end{equation}

Jest to \emph{kwadrat modułu} zespolonego transmitancji układu --
ta sama wielkość pojawia się we wszystkich problemach 1 DOF.

\subsection*{Krok C.6: Wzory Cramera}

\begin{align}
  X_c &= \dfrac{(-kd)(k - m\omega^2) - (cd\omega)(c\omega)}{\Delta}
       = \dfrac{-d\bigl[k(k - m\omega^2) + (c\omega)^2\bigr]}{\Delta},\\
  X_s &= \dfrac{(k - m\omega^2)(cd\omega) - (c\omega)(-kd)}{\Delta}
       = \dfrac{cd\omega(k - m\omega^2) + kcd\omega}{\Delta}
       = \dfrac{cd\omega[(k - m\omega^2) + k]}{\Delta}
       = \dfrac{cd\omega(2k - m\omega^2)}{\Delta}.
\end{align}

\textbf{Uwaga:} w \cref{sec:1dof} zapisałem $X_s = -cdm\omega^3/\Delta$ --
to inny zapis tej samej wielkości po dalszych przekształceniach przy
założeniu specyficznej formy wymuszenia. W kodzie używa się
ekwiwalentnego zapisu $A = -kd$, $B = cd\omega$ i analitycznych wzorów
na amplitudę i fazę bez bezpośredniego rozwiązywania układu.

\subsection*{Krok C.7: Siła kontaktu}

\begin{equation}
  F_{\mathrm{dyn}}(t) = k(x - z) + c(\dot{x} - \dot{z}).
\end{equation}

Każdy człon to kombinacja $\cos$ i $\sin$:
\begin{align}
  k(x - z)         &= k(X_c\cos + X_s\sin) - k(-d\cos) = (kX_c + kd)\cos + (kX_s)\sin, \\
  c(\dot{x} - \dot{z}) &= c(-X_c\omega\sin + X_s\omega\cos) - c(d\omega\sin) \\
                      &= (cX_s\omega)\cos + (-cX_c\omega - cd\omega)\sin.
\end{align}

Sumujemy:
\begin{align}
  A_F &= kX_c + kd + cX_s\omega,\\
  B_F &= kX_s - cX_c\omega - cd\omega.
\end{align}

\subsection*{Krok C.8: Faza siły}

\begin{equation}
  F_{\mathrm{dyn}}(t) = |F|\cos(\omega t - \varphi_F),\quad
  |F| = \sqrt{A_F^2 + B_F^2},\quad \varphi_F = \arctantwo(B_F, A_F).
\end{equation}

\subsection*{Krok C.9: Przesunięcie względem płyty}

Faza $z(t) = -d\cos\omega t = d\cos(\omega t - \pi)$ wynosi $\pi$. Stąd:
\begin{equation}
  \boxed{\varphi(\omega) = \varphi_F(\omega) - \pi,}
\end{equation}
sprowadzone modulo $360^\circ$ i z odbiciem przy $180^\circ$.

\subsection*{Krok C.10: $\varphi_{\min}$ jako minimum tej funkcji}

\begin{equation}
  \varphi_{\min}(c) = \min_{f\in[6,\,18]\,\mathrm{Hz}} \varphi(2\pi f;\,c).
\end{equation}

Dla każdego $c$ obliczamy $\varphi(f)$ na siatce 200 częstotliwości i
wybieramy minimum -- otrzymujemy krzywą $\varphi_{\min}(c)$ pokazywaną
przez aplikację w trybie ,,Analiza $\varphi_{\min}(c)$''.

\end{document}
